% GENERATED FILE. Edit canonical lessons, then run npm run generate:books.
% canonical-source-hash: fnv1a64:dbb66f96bafbbe06
% canonical-lessons: GE-C131-beschaeftigt, GE-C131-frei, GE-C131-geschlossen, GE-C131-ruhig, GE-C131-eng

\chapter{Busy, Free, Closed, Quiet, Narrow}
\label{ch:ge-busy-free-closed-quiet-narrow}
\hlchaptermodality{131}

\begin{chapteropening}
\textbf{By the end of this chapter:} \emph{I can say busy, free, closed, quiet and narrow.}

Complete the last lesson of chapter 131: i can say busy, free, closed, quiet and narrow.
\end{chapteropening}

\section[beschäftigt]{beschäftigt — busy}

\label{lesson:GE-C131-beschaeftigt}



\begin{quote}
Before the new one: say the German for round, then the German for wide.

Say \emph{sechs}, where chs sounds as an x, and then \emph{haben}.
\end{quote}

\subsection*{You'll want to know: beschäftigt}
\textbf{beschäftigt} — \textquotedblleft{}busy\textquotedblright{}.

\begin{grammarlens}[title={What you've built}]
One more word for this chapter.
\end{grammarlens}

\subsection*{Your turn}
\begin{itemize}
\raggedright
  \item \emph{Say it:} \emph{beschäftigt}
  \item \emph{Say it:} \emph{beschäftigt}, once more
  \item \emph{Say it:} say \emph{beschäftigt}
  \item \emph{Recall:} say the German for round, then the German for wide, then say \emph{beschäftigt} again
\end{itemize}

\begin{tcolorbox}[breakable,colback=teal!4,colframe=teal!35!black,title={Before you move on}]
What does beschäftigt mean? (\textquotedblleft{}Busy\textquotedblright{}.) Say it once more.
\end{tcolorbox}
\section[frei]{frei — free}

\label{lesson:GE-C131-frei}



\begin{quote}
Before the new one: say the German for wide, then the German for busy.
\end{quote}

\subsection*{You'll want to know: frei}
\textbf{frei} — \textquotedblleft{}free\textquotedblright{}.

\begin{grammarlens}[title={What you've built}]
One more word for this chapter.
\end{grammarlens}

\subsection*{Your turn}
\begin{itemize}
\raggedright
  \item \emph{Say it:} \emph{frei}
  \item \emph{Say it:} \emph{frei}, once more
  \item \emph{Say it:} say \emph{frei}
  \item \emph{Recall:} say the German for wide, then the German for busy, then say \emph{frei} again
\end{itemize}

\begin{tcolorbox}[breakable,colback=teal!4,colframe=teal!35!black,title={Before you move on}]
What does frei mean? (\textquotedblleft{}Free\textquotedblright{}.) Say it once more.
\end{tcolorbox}
\section[geschlossen]{geschlossen — closed}

\label{lesson:GE-C131-geschlossen}



\begin{quote}
Before the new one: say the German for busy, then the German for free.
\end{quote}

\subsection*{You'll want to know: geschlossen}
\textbf{geschlossen} — \textquotedblleft{}closed\textquotedblright{}.

\begin{grammarlens}[title={What you've built}]
One more word for this chapter.
\end{grammarlens}

\subsection*{Your turn}
\begin{itemize}
\raggedright
  \item \emph{Say it:} \emph{geschlossen}
  \item \emph{Say it:} \emph{geschlossen}, once more
  \item \emph{Say it:} say \emph{geschlossen}
  \item \emph{Recall:} say the German for busy, then the German for free, then say \emph{geschlossen} again
\end{itemize}

\begin{tcolorbox}[breakable,colback=teal!4,colframe=teal!35!black,title={Before you move on}]
What does geschlossen mean? (\textquotedblleft{}Closed\textquotedblright{}.) Say it once more.
\end{tcolorbox}
\section[ruhig]{ruhig — quiet, calm}

\label{lesson:GE-C131-ruhig}



\begin{quote}
Before the new one: say the German for free, then the German for closed.
\end{quote}

\subsection*{You'll want to know: ruhig}
\textbf{ruhig} — \textquotedblleft{}quiet, calm\textquotedblright{}.

\begin{grammarlens}[title={What you've built}]
One more word for this chapter.
\end{grammarlens}

\subsection*{Your turn}
\begin{itemize}
\raggedright
  \item \emph{Say it:} \emph{ruhig}
  \item \emph{Say it:} \emph{ruhig}, once more
  \item \emph{Say it:} say \emph{ruhig}
  \item \emph{Recall:} say the German for free, then the German for closed, then say \emph{ruhig} again
\end{itemize}

\begin{tcolorbox}[breakable,colback=teal!4,colframe=teal!35!black,title={Before you move on}]
What does ruhig mean? (\textquotedblleft{}Quiet, calm\textquotedblright{}.) Say it once more.
\end{tcolorbox}
\section[eng]{eng — narrow; tight}

\label{lesson:GE-C131-eng}



\begin{quote}
Before the new one: say the German for closed, then the German for quiet.
\end{quote}

\subsection*{You'll want to know: eng}
\textbf{eng} — \textquotedblleft{}narrow; tight\textquotedblright{}.

\begin{grammarlens}[title={What you've built}]
One more word for this chapter.
\end{grammarlens}

\subsection*{Your turn}
\begin{itemize}
\raggedright
  \item \emph{Say it:} \emph{eng}
  \item \emph{Say it:} \emph{eng}, once more
  \item \emph{Say it:} say \emph{eng}
  \item \emph{Recall:} say the German for closed, then the German for quiet, then say \emph{eng} again
\end{itemize}

\begin{tcolorbox}[breakable,colback=teal!4,colframe=teal!35!black,title={Before you move on}]
What does eng mean? (\textquotedblleft{}Narrow; tight\textquotedblright{}.) Say it once more.
\end{tcolorbox}
